\begin{lemma}[A failed comparison has a canonical one-round response]
For any relation $r$, if
$\neg\mathsf{PowerRel}(r)(x,y)$, then there is a legal move $x'$ from $x$
such that every legal move $y'$ from $y$ still satisfies
$\neg\mathsf{PowerRel}(r)(x',y')$.
\end{lemma}
\begin{proof}
Proceed by the atom/node tags of $x$ and $y$, using the four reduction
clauses \noderef{PowerRelAtomAtom}, \noderef{PowerRelAtomNode},
\noderef{PowerRelNodeAtom}, and \noderef{PowerRelNodeNode}.  For two atoms,
take the unchanged left atom; the assumed failure is the required conclusion.
For an atom on the left and a node on the right, failure of the existential
clause says that every right child fails, again with the unchanged left atom.
For a node on the left and an atom on the right, failure of the universal
clause supplies a left child whose comparison fails, and the right atom has
only its unchanged move.  Finally, for two nodes, failure of the
universal-existential clause supplies a left child for which every right
child fails.  In each case \noderef{PowerMove} identifies exactly the legal
moves just described, proving the stated uniform response property.
\end{proof}
